\chapter{Phase 1 - Ingestion}

\section{Phase objective}

Convert a raw documentation source into a list of fragments tagged by
candidate endpoint, regardless of whether that source is unstructured prose
like PDF, HTML, or Markdown, or already structured data like a partial
OpenAPI document, a Postman collection, or a HAR capture. This phase
produces the input that Phase 2 will consume, and defines the second data
contract of the pipeline: the \texttt{Chunk} model.

\section{The \texttt{Chunk} model}

\begin{lstlisting}[language=Python, caption={\texttt{src/spekforge/models/chunk.py}}]
from enum import Enum
from pydantic import BaseModel
from spekforge.models.endpoint import EndpointSpec

class SourceType(str, Enum):
    MARKDOWN = "markdown"
    OPENAPI_PARTIAL = "openapi_partial"
    POSTMAN = "postman"
    HAR = "har"

class Chunk(BaseModel):
    id:str
    source_type: SourceType
    source_path: str
    raw_text: str
    candidate_method: str | None = None
    candidate_path: str | None = None
    pre_extracted: EndpointSpec | None = None
\end{lstlisting}

\texttt{Chunk} lives in \texttt{models/}, alongside \texttt{EndpointSpec},
\texttt{ParamSpec}, and \texttt{SchemaSpec}, because it is a contract
between two stages of the pipeline (Ingestion $\rightarrow$ Extraction),
following the same pattern established in Phase 0.

The fields \texttt{candidate\_method} and \texttt{candidate\_path} are typed
as plain \texttt{str | None}, not as the strict \texttt{HTTPMethod} enum
used inside \texttt{EndpointSpec}: at this stage they are only a heuristic
guess extracted from a heading. This avoids forcing the enum too early,
which would break ingestion the moment a heading did not match the
expected pattern exactly.

The \texttt{pre\_extracted} field is the mechanism that allows already
structured sources to skip the LLM call in Phase 2: when this field is
populated, the orchestrator uses it directly; when it is \texttt{None},
\texttt{raw\_text} is sent to the model instead. This keeps
\texttt{to\_intermediate()} returning a uniform \texttt{list[Chunk]}
regardless of which ingestor produced it.

\section{The \texttt{Ingestor} protocol}

\begin{lstlisting}[language=Python, caption={\texttt{src/spekforge/ingest/base.py}}]
from pathlib import Path
from typing import Protocol, runtime_checkable
from spekforge.models.chunk import Chunk

@runtime_checkable
class Ingestor(Protocol):
    def accepts(self, source: Path) -> bool: ...
    def to_intermediate(self, source: Path) -> list[Chunk]: ...
\end{lstlisting}

The purpose of \texttt{accepts()} is to answer whether a given ingestor can
process a source; in the same way, \texttt{router.py} tries each
registered ingestor in order and dispatches to the first one that returns
\texttt{True}. \texttt{@runtime\_checkable} allows \texttt{isinstance()}
checks against this protocol without requiring explicit inheritance: both
ingestors satisfy the protocol, while a plain object does not.

\section{Markdown chunking heuristic}

\begin{lstlisting}[language=Python, caption={\texttt{src/spekforge/ingest/chunking.py}}]
import re
from spekforge.models.chunk import Chunk, SourceType

_METHODS = "GET|POST|PUT|PATCH|DELETE|HEAD|OPTIONS"
_ENDPOINT_HEADING = re.compile(
    rf"^#{{1,6}}\s*\**`?({_METHODS})`?\**\s+`?(/\S*?)`?\s*$",
    re.IGNORECASE | re.MULTILINE,
)

def split_into_chunks(markdown:str, source_path: str) -> list[Chunk]:
    matches = list(_ENDPOINT_HEADING.finditer(markdown))

    chunks: list[Chunk] = []
    for index, match in enumerate(matches):
        start = match.start()
        end = matches[index + 1].start() if index +1 < len(matches) else len(markdown)
        raw_text = markdown[start:end].strip()

        chunks.append(
            Chunk(
                id=f"{source_path}#{index}",
                source_type=SourceType.MARKDOWN,
                source_path=source_path,
                raw_text=raw_text,
                candidate_method=match.group(1).upper(),
                candidate_path=match.group(2).strip(),
            )
        )
    return chunks
\end{lstlisting}

The boundary between chunks is a heading, not any mention of a method and
path elsewhere in the text. The regular expression is anchored with the
pattern \texttt{\^{}...\$} in \texttt{MULTILINE} mode and requires the
entire line to be the heading, which prevents an inline sentence that
happens to mention an HTTP method and a path from triggering a false
boundary inside a paragraph. The pattern also tolerates the decorations
commonly produced by Markdown documentation generators, such as backticks
and bold markers around the method and path.

This design carries a known limitation: if a documentation source does not
follow this heading convention, \texttt{split\_into\_chunks} returns an
empty list. No fallback has been implemented yet; the idea for a future
expansion is to build one, but only once there is real data to validate
what a sensible fallback should look like.

\section{\texttt{MarkitdownIngestor}}

\begin{lstlisting}[language=Python, caption={\texttt{src/spekforge/ingest/markitdown\_ingestor.py}}]
from pathlib import Path
from markitdown import MarkItDown
from spekforge.ingest.chunking import split_into_chunks
from spekforge.models.chunk import Chunk

_SUPPORTED_EXTENSIONS = {".pdf", ".docx", ".pptx", ".html", ".htm", ".md"}

class MarkitdownIngestor:
    def __init__(self) -> None:
        self.converter = MarkItDown()

    def accepts(self, source: Path) -> bool:
        return source.suffix.lower() in _SUPPORTED_EXTENSIONS

    def to_intermediate(self, source: Path) -> list[Chunk]:
        result = self.converter.convert(str(source))
        return split_into_chunks(result.markdown, str(source))
\end{lstlisting}

This ingestor accepts the set of extensions listed in
\texttt{\_SUPPORTED\_EXTENSIONS}, converts the source to Markdown through
\texttt{markitdown}, and delegates the boundary-detection logic to
\texttt{chunking.py}. Extension recognition was verified for all six
supported extensions, and the full conversion path was additionally
validated end to end against real \texttt{.html} and \texttt{.md}
fixtures, confirming that raw input is correctly converted and
boundary-split into chunks; \texttt{.pdf}, \texttt{.docx}, and
\texttt{.pptx} were checked only for extension acceptance so far.

\section{\texttt{StructuredIngestor}}

This ingestor handles three distinct already-structured formats and
converts each one directly into \texttt{EndpointSpec} objects, bypassing
the LLM entirely for that data.

\subsection{Format detection}

\begin{lstlisting}[language=Python,caption={Format detection heuristic}]
def _detect_format(data: dict) -> SourceType | None:
    if not isinstance(data, dict):
        return None
    if "paths" in data and ("openapi" in data or "swagger" in data):
        return SourceType.OPENAPI_PARTIAL
    if "item" in data and "info" in data:
        return SourceType.POSTMAN
    if "entries" in data.get("log", {}):
        return SourceType.HAR
    return None
\end{lstlisting}

Detection relies on top-level JSON keys, because all three formats are
plain \texttt{.json} (or, in the case of HAR, sometimes \texttt{.har}) and
are only distinguishable by their shape.

\subsection{OpenAPI partial}

\begin{lstlisting}[language=Python, caption={Converting OpenAPI into an EndpointSpec}]
def _from_openapi(data: dict) -> list[tuple[EndpointSpec, dict]]:
    results = []
    for path, path_item in data.get("paths", {}).items():
        for method_str, operation in path_item.items():
            if method_str.upper() not in HTTPMethod.__members__:
                continue

            endpoint = EndpointSpec(
                method=HTTPMethod(method_str.lower()),
                path=path,
                summary=operation.get("summary"),
                description=operation.get("description"),
                parameters=[_openapi_param(p) for p in operation.get("parameters", [])],
                request_body=_openapi_request_body(operation.get("requestBody")),
                responses=_openapi_responses(operation.get("responses", {})),
            )
            results.append((endpoint, {path: {method_str: operation}}))
    return results
\end{lstlisting}

A recursive helper, \texttt{\_openapi\_schema}, translates nested JSON
Schema objects (\texttt{properties}, \texttt{items}, \texttt{required})
into \texttt{SchemaSpec} instances, mirroring the recursive structure
already defined in Phase 0. If a response is declared without a
\texttt{content} block, such as a 204, it is correctly omitted from
\texttt{responses} instead of being assigned an invented schema.

\subsection{Postman collections}

Postman collections nest requests inside folders arbitrarily deep, so
traversal is recursive, as shown below:

\begin{lstlisting}[language=Python, caption={Recursive traversal, with validation of malformed items}]
def _walk_postman_items(items: list, results: list[tuple[EndpointSpec, dict]]) -> None:
    for item in items:
        if "item" in item:
            _walk_postman_items(item["item"], results)
            continue
        if "request" not in item:
            continue
        request = item.get("request", {})
        method_str = request.get("method", "GET").upper()
        if method_str not in HTTPMethod.__members__:
            continue

        url = request.get("url", "")
        endpoint = EndpointSpec(
            method=HTTPMethod(method_str.lower()),
            path=_postman_path(url),
            summary=item.get("name"),
            parameters=_postman_params(request, url),
        )
        results.append((endpoint, item))

def _postman_path(url) -> str:
    if isinstance(url, str):
        path = urlparse(url).path
    else:
        path = "/" + "/".join(url.get("path", []))
    segments = [f"{{{p[1:]}}}" if p.startswith(":") else p for p in path.split("/")]
    return "/".join(segments) or "/"

def _postman_params(request: dict, url) -> list[ParamSpec]:
    params = [
        ParamSpec(name=h["key"], location=ParamLocation.HEADER, required=not h.get("disabled", False), type="string")
        for h in request.get("header", [])
        if h.get("key") and not h.get("disabled", False)
    ]
    if isinstance(url, dict):
        params += [
            ParamSpec(name=q["key"], location=ParamLocation.QUERY, required=not q.get("disabled", False), type="string")
            for q in url.get("query", [])
            if q.get("key") and not q.get("disabled", False)
        ]
    return params
\end{lstlisting}

Postman's path-variable syntax is normalized to the OpenAPI style, so that
\texttt{EndpointSpec.path} follows a single convention regardless of the
originating format. Postman collections do not carry real response
schemas, so that field is intentionally left empty for this format. Item
traversal was validated against malformed items and against header/query
parameters explicitly marked as disabled in the collection.

\subsection{HAR archives}

HAR is the only one of the three formats that carries a real, captured
response body, which makes it the most reliable source for inferring
response schemas from actual data rather than declared types:

\begin{lstlisting}[language=Python, caption={Inferring a schema from a captured JSON value}]
def _value_to_schema(value) -> SchemaSpec:
    if isinstance(value, dict):
        return SchemaSpec(
            type="object",
            properties={k: _value_to_schema(v) for k, v in value.items()},
            required=list(value.keys()),
        )
    if isinstance(value, list):
        return SchemaSpec(type="array", items=_value_to_schema(value[0]) if value else None)
    if isinstance(value, bool):
        return SchemaSpec(type="boolean")
    if isinstance(value, int):
        return SchemaSpec(type="integer")
    if isinstance(value, float):
        return SchemaSpec(type="number")
    if value is None:
        return SchemaSpec(type="null")
    return SchemaSpec(type="string")
\end{lstlisting}

This function checks the type of the captured value in sequence, starting
with \texttt{bool} before \texttt{int} so that a JSON boolean is not
misclassified as an integer, and returns the corresponding
\texttt{SchemaSpec}. It was also verified that an entry declaring an
invalid HTTP method is skipped rather than processed.

\section{Router}
\begin{lstlisting}[language=Python, caption={\texttt{src/spekforge/ingest/router.py}}]
from pathlib import Path
from spekforge.ingest.base import Ingestor
from spekforge.ingest.markitdown_ingestor import MarkitdownIngestor
from spekforge.ingest.structured_ingestor import StructuredIngestor
from spekforge.models.chunk import Chunk

_INGESTORS: list[Ingestor] = [StructuredIngestor(), MarkitdownIngestor()]

def route(source: Path) -> list[Chunk]:
    for ingestor in _INGESTORS:
        if ingestor.accepts(source):
            return ingestor.to_intermediate(source)

    raise ValueError(f"No ingestor accepts source: {source}")
\end{lstlisting}

\texttt{route()} is in charge of selecting the ingestor depending on the
source. \texttt{StructuredIngestor} is tried first, since it is the more
specific of the two, validating the actual JSON content shape rather than
just the extension. If no registered ingestor accepts the source, the
function raises a \texttt{ValueError} including the source path. This also
covers a source path that does not exist on disk: \texttt{accepts()}
catches the \texttt{OSError} raised while attempting to read it and simply
returns \texttt{False}, so a missing file surfaces as the same explicit
error as any other unrecognized source.

\section{Validation against real documentation}

To validate the pipeline against something more realistic than hand-written
fixtures, the official OpenAPI 3.0 specification for Swagger Petstore was
retrieved directly from its public instance
(\texttt{petstore3.swagger.io}) and stored as the hidden ground truth for
this project, in \texttt{tests/fixtures/petstore/ground\_truth/}.

The Swagger UI documentation portal itself, however, could not be used
directly as an ingestion source: it is a single-page application whose
static HTML contains nothing but an empty \texttt{div} placeholder, with
all real content rendered client-side by JavaScript. To still validate the
pipeline against
realistic narrative documentation, a Markdown document was generated
directly from the official specification's 19 operations, following the
same heading-per-endpoint convention real documentation generators
produce, and stored separately in
\texttt{tests/fixtures/petstore/input\_docs/}, apart from the hidden
ground truth.

Running the full pipeline (\texttt{router.route()} $\rightarrow$
\texttt{MarkitdownIngestor} $\rightarrow$ \texttt{chunking.py}) against
this generated Markdown fixture detected all 19 operations defined in the
official specification, with zero missed and zero spurious endpoints. The
next figure shows this ingestion run.

\realcapture{0.85}{Phase1Router.png}{Full ingestion pipeline output against the real Petstore documentation fixture: 19 of 19 endpoints detected}

\subsection{Edge cases review}

Beyond the fixture above, every component of this phase was additionally
exercised against edge cases designed to probe the phase further, such as
empty input; boundary conditions in heading detection like consecutive
headings, duplicate headings, and unusual whitespace; invalid methods;
extension case sensitivity; malformed and non-existing source files;
OpenAPI responses without a body; and invalid HTTP methods in Postman and
HAR entries. This review is what led to the guard clauses already shown in
this chapter's code, such as the malformed-item check in
\texttt{\_walk\_postman\_items} and the \texttt{disabled} filter in
\texttt{\_postman\_params}; with those in place, every case listed above
is handled correctly.

\section{Tests}
\begin{lstlisting}[language=Python, caption={\texttt{tests/unit/test\_ingest.py}}]
import json
from pathlib import Path

import pytest
from spekforge.ingest.chunking import split_into_chunks
from spekforge.ingest.markitdown_ingestor import MarkitdownIngestor
from spekforge.ingest.router import route
from spekforge.ingest.structured_ingestor import StructuredIngestor
from spekforge.models.chunk import SourceType

def test_split_into_chunks_detects_endpoints_by_heading():
    markdown = """## GET /pets/{id}
Retrieve a pet by id.

## POST /pets
Create a new pet.
"""
    chunks = split_into_chunks(markdown, "petstore.md")

    assert len(chunks) == 2
    assert chunks[0].candidate_method == "GET"
    assert chunks[0].candidate_path == "/pets/{id}"
    assert chunks[1].candidate_method == "POST"
    assert chunks[1].candidate_path == "/pets"

def test_split_into_chunks_ignores_inline_mentions():
    markdown = "You can call GET /pets/{id} to retrieve a pet"
    chunks = split_into_chunks(markdown, "petstore.md")

    assert chunks == []

def test_structured_ingestor_rejects_generic_json(tmp_path):
    source = tmp_path / "not_structured.json"
    source.write_text(json.dumps({"hello": "world"}))

    assert StructuredIngestor().accepts(source) is False
\end{lstlisting}

In total, nine tests cover this phase's contract, including heading-boundary
detection and its negative case, extension acceptance, all three
\texttt{StructuredIngestor} formats including HAR type inference, generic
JSON rejection, and the router's dispatch and error paths. Together with
the edge-case review, this test suite is what gives confidence that the
ingestion of real documentation behaves as intended.
